\section{LIS3DH Code}
\label{ann:lis3dh_code}

\begin{lstlisting}[language=C, label={code:LIS3DHState}, caption={LIS3DH device state structure}]
typedef struct LIS3DHState 
{
    /* Parent object */
    I2CSlave i2c;

    lis3dh_config_t *config;    /**< Operating mode, scale, ODR */
    lis3dh_params_t *params;    /**< Sensitivity, range, shifts */
    lis3dh_i2c_params_t i2c;    /**< Address pointer, auto-inc  */
    
    /* GPIO outputs */
    qemu_irq int1;               /**< INT1 output line */
    qemu_irq int2;               /**< INT2 output line */

    /* Registers */
    uint8_t status_reg_aux; // Status Register
    uint8_t adc_1_l;    // 1-Axis Acceleration Data Low Register
    uint8_t adc_1_h;    // 1-Axis Acceleration Data High Register
    uint8_t adc_2_l;    // 2-Axis Acceleration Data Low Register
    uint8_t adc_2_h;    // 2-Axis Acceleration Data High Register
    uint8_t adc_3_l;    // 3-Axis Acceleration Data Low Register
    uint8_t adc_3_h;    // 3-Axis Acceleration Data High Register
    :
    :
    :
    :
    uint8_t act_ths;
    uint8_t act_dur;    
    
} LIS3DHState;
\end{lstlisting}


\begin{lstlisting}[language=C, label={code:lis3dh_i2c_event},caption={LIS3DH I2C slave event handler}]
static int lis3dh_i2c_event(I2CSlave *i2c, enum i2c_event event)
{
    LIS3DHState *lis3dh = LIS3DH(i2c);
    
    switch (event) 
    {
        case I2C_START_SEND:    // Start of write operation
            /* Master is starting a WRITE operation 
            (sending data to the device) */
            lis3dh->i2c_params.ptr             = 0xFF;
			lis3dh->i2c_params.auto_increment  = false;
			lis3dh->i2c_params.address_phase   = true;
            break;
            
        case I2C_START_RECV:    // Start of read operation
            /* Master is starting a READ operation 
            (requesting data from the device) */
			if (lis3dh->i2c_params.ptr == 0xFF)
				lis3dh->i2c_params.ptr = LIS3DH_ADDR_WHO_AM_I;

			lis3dh->i2c_params.address_phase = false;
            break;
            
        case I2C_FINISH:        // Stop condition
            /* Master ends the transaction 
            (STOP condition) */
            break;
            
        case I2C_NACK:          // NACK received
            /* Master didn't acknowledge the data */
            break;
        
        default:               // Should never happen
            return -1;
    }

    return 0;
}
\end{lstlisting}

\begin{lstlisting}[language=C, label={code:lis3dh_recv} caption={LIS3DH I2C slave recv (write from master)}]
static uint8_t lis3dh_i2c_recv(I2CSlave *i2c)
{
    LIS3DHState *lis3dh = LIS3DH(i2c);

    uint8_t value = lis3dh_read_register( lis3dh );

    if (lis3dh->i2c_params.auto_increment)
    {
        lis3dh->i2c_params.ptr++;
    }

	return value;
}
\end{lstlisting}

\begin{lstlisting}[language=C, label={code:lis3dh_i2c_send} caption={LIS3DH I2C slave send (read by master)}]
static int lis3dh_i2c_send(I2CSlave *i2c, uint8_t data)
{
    LIS3DHState *lis3dh = LIS3DH(i2c);

	if (lis3dh->i2c_params.ptr == 0xFF && lis3dh->i2c_params.address_phase)
	{
		lis3dh->i2c_params.ptr = data & LIS3DH_SUB_REG_MASK;
		lis3dh->i2c_params.auto_increment =  (data & LIS3DH_SUB_AUTO_INC_MASK) != 0x00 ? true : false;
		lis3dh->i2c_params.address_phase = false;
	}else
	{
		lis3dh_write_register( lis3dh, lis3dh->i2c_params.ptr, data);
		if (lis3dh->i2c_params.auto_increment)
			lis3dh->i2c_params.ptr++;
	}
	return 0;
}
\end{lstlisting}

\begin{lstlisting}[language=C, label={code:lis3dh_qom_send} caption={LIS3DH QOM properties}]
static void lis3dh_initfn(Object *obj)
{
    object_property_add(obj, "accel-x", "int",
                        lis3dh_get_accel_x,
                        lis3dh_set_accel_x, NULL, NULL);
    object_property_add(obj, "accel-y", "int",
                        lis3dh_get_accel_y, 
                        lis3dh_set_accel_y, NULL, NULL);
    object_property_add(obj, "accel-z", "int",
                        lis3dh_get_accel_z,
                        lis3dh_set_accel_z, NULL, NULL);
    object_property_add(obj, "temp", "int",
                        lis3dh_get_temp,
                        lis3dh_set_temp, NULL, NULL);
}
\end{lstlisting}